\section{Introduction}
This document provides a comprehensive overview of the \textbf{Pole Anomaly Project}, an anomaly detection and classification system for nuclear gamma dosimetry signals. The overall objective is to automatically identify and classify transient anomalous events---such as radiation spikes with distinct geometric signatures---embedded within continuous, noisy gamma dose rate measurements.

The project follows a two-phase approach. First, a \textbf{synthetic data generation pipeline} produces large-scale training datasets by injecting user-defined anomaly templates (Bell curves, Squares, M-shapes, Exponential ramps, etc.) into realistic simulated radiation backgrounds derived from real-world sensor baselines. Second, multiple \textbf{1D Deep Learning architectures (1D-CNN, ResNet, CNN-Attention)} are trained on these synthetic datasets to perform multi-class anomaly classification in a streaming inference setting.

The repository is organized into three main components:
\begin{itemize}
    \item \texttt{data-gen/}: The synthetic data generation pipeline, including anomaly templates, the stochastic generator engine, and visualization scripts.
    \item \texttt{architectures/}: The neural network implementations, including the model architectures (1D-CNN, ResNet, CNN-Attention), training loop, inference engine, and robustness testing scripts.
    \item \texttt{documentation/}: This document and associated media.
\end{itemize}
